% chapters/ch00_introduction.tex — Front matter: welcome, how to use the book,
% the readiness ladder, and the 12-week study plan (PLAN.md Section 4, printed as
% the book's introduction). Unnumbered (\chapter*) so Chapter 1 stays Chapter 1.
\chapter*{Start Here --- How to Use This Book}
\addcontentsline{toc}{chapter}{Start Here --- How to Use This Book}
\markboth{Start Here}{Start Here}

\begin{keyidea}
You don't need to be a genius to walk into an embedded/avionics interview and
\emph{own} it. You need the right answers rehearsed until they're reflexes, the
right code in your fingers, and your own projects scripted so you can defend
every line. That's the entire job of this book --- and it's built around
\emph{your} work: DFCC, CCDL, SPIL, RS422 Embedded C, the BEL DPRAM and 1553B
automation.
\end{keyidea}

\noindent Welcome. This is a \textbf{do-it-don't-just-read-it} book. If you read
it like a novel you'll forget it; if you \emph{work} it --- speak the answers
aloud, attempt every coding problem on paper, take the scorecards honestly ---
you'll move from ``I know roughly how this works'' to ``I can explain it to a
principal engineer and write the code on the whiteboard.'' Twelve weeks, and the
needle goes from \textbf{10\% ready} to \textbf{100\% ready}.

\section*{The promise (and the deal)}
Every chapter is engineered to produce two things on every page: an
\emph{interview answer} or \emph{working code}. No history lessons, no filler.
In return, you commit to three non-negotiable habits:

\begin{enumerate}
\item \textbf{Attempt before reading.} For every coding problem, cover the
  solution and work it on paper for at least 10 minutes first. Reading solutions
  cold feels productive and teaches you almost nothing.
\item \textbf{Speak it out loud.} Interviews are spoken, not read. Say each
  Tier-2 answer aloud, once, as if to a panel.
\item \textbf{8 out of 10 to advance.} Each chapter ends with a 10-question
  scorecard. Hit 8/10 before moving on; revisit a failed scorecard after three
  days (spaced repetition).
\end{enumerate}

\section*{How every chapter is built}
Once you know the rhythm, the book is easy to navigate. Each teaching chapter
(1--12, 14) has the same five parts, always in this order:

\begin{description}
\item[\textcolor{accent}{Primer}] The minimum theory to answer the questions ---
  a few pages, no more. Read it once to load the concepts.
\item[\textcolor{accent}{Q\&A Bank, in three tiers}] The heart of the book.
  \emph{Tier~1} = screening / phone-round questions (answer in 2--4 sentences).
  \emph{Tier~2} = core panel-round questions (a short paragraph). \emph{Tier~3}
  = principal-engineer depth, each ending with \textbf{``the trap''} it defends
  against. Work up the tiers; if Tier~3 feels hard, that's the point.
\item[\textcolor{accent}{Coding Gym}] Problems in the format \emph{Problem
  $\rightarrow$ Hints $\rightarrow$} \textsc{attempt before reading on}
  \emph{$\rightarrow$ Solution with commentary $\rightarrow$ interviewer
  follow-up variants.} Every solution in this book compiles clean under
  \code{gcc -Wall -Wextra -std=c11} (the Python under \code{python3}) --- it was
  run before it was printed.
\item[\textcolor{okgreen!70!black}{Link to Her Resume}] A green box with
  ready-to-say sentences connecting the topic to your actual DFCC/CCDL/SPIL/BEL
  work --- so theory becomes \emph{your} story.
\item[\textcolor{accent}{Self-Test Scorecard}] Ten rapid-fire questions and an
  answer key. Your 8/10 gate.
\end{description}

\noindent Two special chapters break the pattern on purpose: \textbf{Chapter 13}
is a 60-problem Coding Gym (pure drill), and \textbf{Chapter 17} is five full
mock interviews (run them like the real thing). And a friendly heads-up:
\textbf{Chapters 15 \& 16} (your project deep-dives and behavioural answers) are
deliberately left as fill-in skeletons --- they must be built from \emph{your}
real facts in \code{inputs/INTAKE.md}, never invented, because a made-up detail
is the fastest way to get caught. Fill that file in and those two chapters
complete themselves.

\section*{The readiness ladder}
\begin{center}
\small
\begin{tabular}{@{}rl@{}}
\toprule
\textbf{Where you are} & \textbf{What it feels like}\\
\midrule
10\% & ``I've used this but couldn't teach it.''\\
35\% & ``I can explain the C fundamentals and the memory model.''\\
50\% & ``Serial, SPI, I2C --- I can draw the frames and do the baud math.''\\
65\% & ``CAN, 1553/429, RTOS --- I can reason about timing and redundancy.''\\
75\% & ``Control, debugging, V\&V --- I can defend a process, not just a fact.''\\
88\% & ``I can write the ring buffer and parse a frame under pressure.''\\
95\% & ``My projects are scripted; I know where my knowledge ends and say so.''\\
100\% & ``Bring on the panel.''\\
\bottomrule
\end{tabular}
\end{center}

\section*{The 12-week study plan}
About 22 hours a week (\,$\approx$2.5\,h on weekdays, 5\,h on weekend days\,).
Follow the order --- each part builds on the last.

\begin{center}
\small
\begin{tabular}{@{}clll@{}}
\toprule
\textbf{Weeks} & \textbf{Focus} & \textbf{Chapters} & \textbf{Readiness}\\
\midrule
1--3  & Embedded C, memory, ARM            & 1, 2, 3      & 10\% $\rightarrow$ 35\%\\
4--5  & Peripherals \& serial protocols    & 4, 5, 6      & $\rightarrow$ 50\%\\
6--7  & CAN, avionics buses, RTOS          & 7, 8, 9      & $\rightarrow$ 65\%\\
8     & Control, debugging, V\&V           & 10, 11, 12   & $\rightarrow$ 75\%\\
9--10 & Coding Gym daily + Python          & 13, 14       & $\rightarrow$ 88\%\\
11    & Project \& behavioural narratives  & 15, 16       & $\rightarrow$ 95\%\\
12    & All five mock interviews           & 17           & $\rightarrow$ 100\%\\
\bottomrule
\end{tabular}
\end{center}

\noindent From week 9, re-test one old scorecard each day as a warm-up --- that
spaced repetition is what makes the knowledge \emph{stick} under pressure.

\section*{A word before you start}
This is a lot of material, and that's normal --- nobody absorbs it in a weekend,
and you're not supposed to. Go part by part, tick the scorecards, and trust the
ladder. On a hard day, do one Coding Gym problem and one Tier-2 answer out loud;
momentum beats marathon. You already do this work for real --- this book just
teaches you to \emph{say} it the way interviewers need to hear it. Turn the
page and let's get you to 100\%.

\begin{flushright}
\itshape --- Your interview-prep cornerman
\end{flushright}
